%=== ABSTRACT ===
\newpage
\chapter*{\centering Abstract}
\markboth{Abstract}{}
\addcontentsline{toc}{chapter}{Abstract}

\begin{spacing}{1.5}
\setlength{\parskip}{0.18in}

Delivery quadrotors experience abrupt and coupled changes in mass, inertia and
centre of mass when a payload is grasped or released. A nonlinear model
predictive controller (NMPC) that retains the unloaded model applies an incorrect
hover thrust, underestimates the moment demanded by an eccentric load, and loses
the intended constraint margin. This dissertation develops a
moving-horizon-estimation-enabled NMPC architecture that adapts the predictive
model online from ordinary flight signals alone, and --- the central claim ---
that never treats an attach or release command as evidence that the payload has
changed.

The controller issues the gripper release command itself and still does not use
it to clear its payload model. The motivation is that a release command is a
request while the model must describe the vehicle: a gripper can be slow, can
jam, can release part of a load, and a load can be lost with no command at all.
A command-triggered controller then flies a loaded vehicle with an empty model,
silently.

Making that refusal affordable requires a payload parameterisation that remains
well posed through release. The conventional pair --- payload mass and grasp
offset --- is singular in the limit that matters: as the payload mass tends to
zero the offset becomes undefined, and an estimator holding it fixed can cancel
the resulting phantom trim moment only by corrupting its mass estimate, which was
observed as 14 of 55 post-release samples resting on the mass lower bound. The
estimator developed here augments the composite mass together with the horizontal
\emph{first mass moment} $s=m_P r_{xy}$, the product that the equations of motion
actually contain, giving a 16-state moving horizon estimator over a
\SI{2.0}{\second} window at \SI{10}{\hertz}. The centre-of-mass offset then
follows as $c=s/m$ with no division by the payload mass and no knowledge of the
grasp geometry; the off-diagonal inertia terms follow with the payload mass
cancelling exactly; and the dominant roll/pitch increment follows from the
composite mass and a weak vertical-arm prior. Two reductions are derived and
quantified: the second-order horizontal inertia terms carry a factor $1/m_P$ in
these coordinates and are discarded for a \SI{2.6}{\percent} bias, and the
remaining mass-dependent inertia coefficient, which otherwise lets the optimiser
dilute a phantom moment by inflating the mass, is frozen within each window.

The physical wrench driving the estimator --- collective thrust, roll and pitch
moments, and the yaw reaction --- is reconstructed from rotor speeds, so the
estimator is never driven by a signal computed from its own output. A window
vertical-force balance supplies a prior-free re-initialisation that uses only the
two velocity endpoints of the window. One atomic, health-gated estimate frame
feeds both the 13-state NMPC (20 stages at \SI{0.05}{\second}, \SI{20}{\hertz})
and a continuous body-rate command scaling with explicit headroom limiting. The
lifecycle is split between two state machines: the estimator owns payload
presence and may declare release on a collapse of the first moment relative to a
self-formed, frozen reference; the controller owns the model and clears it only
on persistent confidence. An unconfirmed release becomes \textsc{unresolved} --- a
hold with a smooth brake to hover --- never an assumption of completion.

The system is implemented in ROS~2 around PX4 software-in-the-loop, Gazebo and
acados. In the frozen campaign the interface confirmed release in 32 of 36
eligible flights through three confirmation paths, with median delays of
\SI{3.46}{\second} in manoeuvre and \SI{16.56}{\second} after braking to hover.
Under a \SI{4}{\second} delayed detachment a command-triggered baseline cleared
its model early in 5 of 5 pairs while the proposed interface did so in none of
six runs. Against a variant armed only by the command, eventlessness detected 10
of 18 unsignaled payload losses that the command-armed variant cannot see, with
no false declaration in 209 loaded flights. The costs are reported as plainly as
the benefits: confirmation is delayed and the delay is bimodal, roughly
\SI{40}{\percent} of flights never form the moment reference the detector needs,
and a \SI{\pm5}{\percent} error in the estimator's thrust coefficient defeats
payload recognition in both directions.

\textbf{Keywords:} quadrotor; nonlinear model predictive control; moving horizon
estimation; first mass moment; inertial parameter identification; eventless
payload adaptation; autonomous release detection; acados; PX4 SITL.
\end{spacing}
\newpage
%=== END OF ABSTRACT ===
